\chapter{Thực nghiệm và phân tích kết quả}\label{chap:evaluation}

\section{Thiết kế thực nghiệm}
Báo cáo kiểm tra ba câu hỏi: hệ thống có xuất được assessment và proof
cho dữ liệu mẫu hay không; năm luật bao phủ bao nhiêu cặp nhãn gold;
và độ chính xác trong phần được đánh giá là bao nhiêu.
Các kết quả dưới đây được chạy lại từ source ngày 07/09/2026
bằng Python 3.12.14 trong môi trường \src{.venv}.

\begin{table}[htbp]
\centering\small
\caption{Các dữ liệu và kiểm tra được sử dụng}\label{tab:data}
\begin{tabularx}{\textwidth}{p{3.3cm}p{2.0cm}Y}
\toprule
Thành phần & Quy mô & Mục đích và giới hạn\\
\midrule
Fixture JSON & 5 mapping & Minh họa chức năng build, query, explain và bốn mức đầu ra.\\
Gold XLSX & 500 cặp nhãn & Đo khả năng phân loại A/B/C; không đo khả năng tìm mã đích.\\
Unit tests & 19 test & Kiểm tra các nhánh code; SHACL là một test tùy chọn.\\
Graph store & 0 quad & Trạng thái đọc tại lần kiểm tra; chưa cung cấp dữ liệu thực nghiệm đầy đủ.\\
TTL nguồn & LFS pointer & Khai báo object 262.786.162 byte, chưa phải nội dung RDF đầy đủ.\\
\bottomrule
\end{tabularx}
\end{table}

Gold nằm tại
\src{gold_datasets/mapping_level/code_pair_mapping_level__gold_standard__reasoning.xlsx}.
Evaluator đọc hai nhãn và nhãn chuẩn từ sheet đầu tiên, dựng CodeConcept
với mã giả S0/T0, S1/T1, rồi gọi reasoner.
Trường gold chỉ dùng để so sánh sau dự đoán, không được đưa vào điều kiện luật.
Tên evidence \texttt{mapping\_level\_gold\_standard} chỉ là metadata của
fixture đánh giá. Tập dữ liệu này là bộ tham chiếu sẵn có; không có thiết kế
chia train/validation/test mới hay khảo sát chuyên gia độc lập trong đồ án.

\section{Các chỉ số}
Gọi $N$ là số cặp, $g_i\in\{A,B,C\}$ là nhãn chuẩn,
$\hat y_i\in\{A,B,C,\bot\}$ là dự đoán.
Đặt $a_i=\Ind[\hat y_i\ne\bot]$ và
$c_i=\Ind[\hat y_i=g_i]$. Khi đó
\begin{align}
 \operatorname{Coverage} &= \frac{\sum_i a_i}{N},\\
 \operatorname{Accuracy}_{\mathrm{all}} &= \frac{\sum_i c_i}{N},\\
 \operatorname{Accuracy}_{\mathrm{assessed}}
 &= \frac{\sum_i c_i}{\sum_i a_i},\\
 \operatorname{Abstention} &= 1-\operatorname{Coverage}.
\end{align}
Trong code, tỷ số có mẫu số bằng 0 được trả về 0.
Với tập gold chỉ gồm A/B/C, có đẳng thức
\begin{equation}
 \operatorname{Accuracy}_{\mathrm{all}}
 =\operatorname{Coverage}\cdot
  \operatorname{Accuracy}_{\mathrm{assessed}}.
 \label{eq:metricrelation}
\end{equation}
Việc coi \Unassessed{} là không đúng trong accuracy toàn tập không
biến nó thành một lớp bệnh hay một nhãn gold mới.

Với từng mức $\ell\in\{A,B,C\}$,
\begin{equation}
 P_\ell=\frac{TP_\ell}{TP_\ell+FP_\ell},\quad
 R_\ell=\frac{TP_\ell}{TP_\ell+FN_\ell},\quad
 F_{1,\ell}=\frac{2P_\ell R_\ell}{P_\ell+R_\ell}.
\end{equation}
Các cặp gold mức $\ell$ nhưng bị từ chối vẫn nằm trong $FN_\ell$.
Do đó recall phản ánh cả giới hạn bao phủ.

\section{Kết quả chức năng trên năm mapping}
\begin{table}[htbp]
\centering\small
\caption{Đầu ra tái lập trên fixture có sẵn}\label{tab:sample}
\begin{tabularx}{\textwidth}{lllY}
\toprule
Nguồn & Đích & Đầu ra & Luật được chọn\\
\midrule
K05.1 & DA0B.Y & A & Exact label.\\
K05.1 & DA0Z & $\bot$ & Insufficient evidence.\\
K25.9 & DA60.Y & C & Explicit qualifier conflict.\\
A31 & 1B21.Z & A & Exact label.\\
A31 & 1H0Z & B & Generalized condition.\\
\bottomrule
\end{tabularx}
\end{table}
Kiểm tra Python của fixture trả về conforms=true, không có error và warning.
Năm kết quả có đủ conclusion, premise, rule ID và source trong cấu trúc proof.
Đây là kiểm tra chức năng trên ví dụ minh họa; không diễn giải thành
độ chính xác 100\% trên một bộ kiểm thử y sinh độc lập.

Với K05.1 đến DA0Z, hai tập từ không có phần giao phù hợp để kích hoạt
luật; hệ thống trả $\bot$ dù nhãn đích chứa “unspecified”.
Điều này minh họa rằng có tên mã và có từ chỉ tính khái quát chưa đủ
để reasoner suy ra mức B.

Trong ví dụ K25.9 đến DA60.Y, cụm “without hemorrhage” xuất hiện
ở nhãn nguồn; nhãn đích có từ “haemorrhagic”.
Phép chuẩn hóa và tìm cụm từ kích hoạt luật xung đột.
Tuy nhiên, luật cũng bắt gặp “chronic” trong
“unspecified as acute or chronic” và có thể đưa vào danh sách
acute/chronic. Trường hợp này cho thấy một đầu ra C có thể đi cùng
một premise phụ chưa được diễn giải ngữ nghĩa chính xác.

\section{Kết quả trên bộ gold}
\begin{table}[htbp]
\centering
\caption{Kết quả tổng hợp của bộ suy luận xác định}\label{tab:metrics}
\begin{tabular}{lrr}
\toprule
Chỉ số & Phân số hoặc số lượng & Giá trị\\
\midrule
Tổng số cặp & 500 & 100,00\%\\
Được đánh giá & 22/500 & 4,40\%\\
Đúng trên toàn bộ tập & 21/500 & 4,20\%\\
Đúng trong phần được đánh giá & 21/22 & 95,45\%\\
Sai trong phần được đánh giá & 1/22 & 4,55\%\\
Chưa được đánh giá & 478/500 & 95,60\%\\
\bottomrule
\end{tabular}
\end{table}

\begin{table}[htbp]
\centering
\caption{Ma trận nhầm lẫn: hàng là gold, cột là dự đoán}\label{tab:confusion}
\begin{tabular}{lrrrrr}
\toprule
Gold & A & B & C & $\bot$ & Tổng\\
\midrule
A & 1 & 1 & 0 & 106 & 108\\
B & 0 & 9 & 0 & 163 & 172\\
C & 0 & 0 & 11 & 209 & 220\\
\midrule
Tổng & 1 & 10 & 11 & 478 & 500\\
\bottomrule
\end{tabular}
\end{table}

\begin{table}[htbp]
\centering
\caption{Precision, recall và F1 theo từng mức; đơn vị \%}\label{tab:perlevel}
\begin{tabular}{lrrr}
\toprule
Mức & Precision & Recall & F1\\
\midrule
A & 100,00 & 0,93 & 1,83\\
B & 90,00 & 5,23 & 9,89\\
C & 100,00 & 5,00 & 9,52\\
\bottomrule
\end{tabular}
\end{table}

\begin{figure}[htbp]
\centering
\begin{tikzpicture}
\begin{axis}[
  xbar stacked,width=0.91\textwidth,height=6cm,
  xmin=0,xmax=240,xtick={0,50,100,150,200},
  xlabel={Số cặp nhãn},ytick={1,2,3},yticklabels={Gold A,Gold B,Gold C},
  ymin=0.4,ymax=3.6,bar width=17pt,
  legend style={at={(0.5,-0.42)},anchor=north,legend columns=3,
    font=\small,draw=none,column sep=6pt},
  tick label style={font=\small},grid=major,grid style={gray!15},
  reverse legend=false]
\addplot[fill=okcolor,draw=none] coordinates {(1,1) (9,2) (11,3)};
\addplot[fill=warncolor,draw=none] coordinates {(1,1) (0,2) (0,3)};
\addplot[fill=gray!40,draw=none] coordinates {(106,1) (163,2) (209,3)};
\legend{Đúng,Sai,Chưa đánh giá}
\end{axis}
\end{tikzpicture}
\caption{Phân bố kết quả theo nhãn chuẩn. Phần chưa đánh giá chiếm ưu thế
ở cả ba nhóm; các số đếm chi tiết nằm trong Bảng~\ref{tab:confusion}.}
\label{fig:coverage}
\end{figure}
\FloatBarrier

Các bảng và Hình~\ref{fig:coverage} thể hiện hai góc nhìn khác nhau:
phần được xử lý có ít lỗi, nhưng phần lớn dữ liệu chưa được xử lý.
Con số 95,45\% chỉ dựa trên 22 cặp; không thể dùng riêng nó để kết luận
hệ thống đạt hiệu quả cao trên toàn bộ miền bài toán.
Precision 100\% của A cũng chỉ dựa trên một dự đoán A.

\section{Mức đóng góp của từng luật}
\begin{table}[htbp]
\centering\small
\caption{Số lần từng luật được chọn trên 500 cặp gold}\label{tab:rulecount}
\begin{tabular}{lrr}
\toprule
Luật & Số lần được chọn & Số dự đoán đúng A/B/C\\
\midrule
Exact label & 1 & 1\\
Explicit qualifier conflict & 11 & 11\\
Compatible specificity & 10 & 9\\
Generalized condition & 0 & 0\\
Insufficient evidence & 478 & Không áp dụng\\
\bottomrule
\end{tabular}
\end{table}
Các số này thống kê \emph{luật được chọn đầu tiên}, không phải số lần
điều kiện của mỗi luật đúng nếu kiểm tra độc lập. Luật generalized có
ví dụ dương trong fixture nhưng không được chọn trên bộ gold.
Đồ án chưa thực hiện ablation thay đổi thứ tự hoặc bỏ từng luật,
nên chưa kết luận được đóng góp nhân quả của mỗi luật tới chất lượng.

\section{Phân tích lỗi cụ thể}
Cặp gold có trường \texttt{Label=378} là:
\begin{samepage}\begin{quote}
Nguồn: Diabetes mellitus type 2.\\
Đích: Type 2 diabetes.\\
Gold: A. Dự đoán: B.
\end{quote}\end{samepage}
Hai tập từ sau chuẩn hóa là
\begin{equation}
 S=\{\text{diabetes,mellitus,type,2}\},\qquad
 T=\{\text{type,2,diabetes}\},\qquad T\subsetneq S.
\end{equation}
Chuỗi không giống hệt nhau nên luật 1 không khớp; không có xung đột
nên luật 3 trả B. Theo gold, khác biệt từ “mellitus” ở cặp này không
làm thay đổi mức tương đương. Lỗi phản ánh khoảng cách giữa bao hàm
từ vựng và quan hệ ngữ nghĩa.

Hướng cải thiện là thêm tri thức đồng nghĩa theo miền, có kiểm thử âm
và đánh giá lại độ bao phủ. Không nên xóa một từ trên mọi nhãn chỉ để
sửa đúng một dòng gold, vì có thể tạo ra lỗi ở những trường hợp chưa quan sát.

\section{Kiểm thử và tính tái lập}
Lần chạy \src{python -m unittest discover -s tests -v} thu được
19 test, gồm 18 đạt và 1 bị bỏ qua do thiếu dependency SHACL tùy chọn.
Các nhóm đã chạy gồm operators, năm nhánh reasoner, repository,
validation, importer với dữ liệu giả lập, XLSX, evaluation và SPARQL guard.
Không báo cáo rằng SHACL đã chạy thành công trong môi trường này.

\placeholderfigure{images/test_results.png}
{Minh họa kết quả kiểm thử và đánh giá tái lập.}
{Chụp terminal hiển thị tổng số unit test, test bị skip và JSON metrics
từ cokb\_eval. Giữ rõ coverage=0.044, overall\_accuracy=0.042,
accuracy\_when\_assessed=0.954545 và unassessed\_count=478.
Nếu cài SHACL và chạy lại, cập nhật trạng thái test trong báo cáo.}
{fig:tests}

Script kèm báo cáo lưu JSON đánh giá, phiên bản Python, hash các file
source liên quan, hash dữ liệu gold, kết quả unit test và artifact mẫu.
Các số trong bảng được đối chiếu với những artifact này.
Không có số đo latency, chi phí LLM, hiệu quả người dùng hoặc kết quả
triển khai trên toàn bộ ICD trong lần thực nghiệm này.
